\documentclass[10pt,conference]{IEEEtran}
\usepackage{amsmath,amssymb,amsfonts}
\usepackage{graphicx}
\usepackage{xcolor}
\usepackage{hyperref}
\usepackage{booktabs}
\usepackage{multirow}
\usepackage{subcaption}
\usepackage{algorithm}
\usepackage{algorithmic}
\usepackage{balance}
\usepackage{tikz}
\usepackage{pgfplots}
\pgfplotsset{compat=1.18}

% Custom commands
\newcommand{\todo}[1]{{\color{red} TODO: #1}}
\newcommand{\fixme}[1]{{\color{blue} FIX: #1}}

% Title information
\title{Paper Title}
\author{\IEEEauthorblockN{Author One}
\IEEEauthorblockA{Department of Computer Science\\University Name\\City, Country\\email@university.edu}
\and
\IEEEauthorblockN{Author Two}
\IEEEauthorblockA{Department of Computer Science\\University Name\\City, Country\\email@university.edu}
}

\begin{document}

\maketitle

\begin{abstract}
This template provides a structured starting point for academic papers.
Replace this text with your abstract (150-250 words).
The abstract should summarize the problem, approach, key results, and significance.
\end{abstract}

\begin{keywords}
Keyword one, keyword two, keyword three
\end{keywords}

\section{Introduction}
\label{sec:introduction}

% Hook: Start with a compelling statement about the problem
The increasing complexity of [problem domain] presents significant challenges for [target audience].

% Problem statement
Despite advances in [related area], existing approaches struggle with [specific limitation].
This limitation is particularly problematic because [reason].

% Research gap
While prior work has addressed [related problem A] \cite{refA} and [related problem B] \cite{refB},
no existing solution effectively handles [specific unsolved problem].

% Our contribution
In this paper, we present [solution name], a novel approach that [brief description of approach].
Our key contributions are:
\begin{itemize}
    \item The first [contribution 1], which enables [benefit].
    \item A novel [contribution 2] that improves [metric] by [percentage].
    \item An efficient [contribution 3] with [complexity] time complexity.
\end{itemize}

% Significance
Our approach is significant because [reason 1], [reason 2], and [reason 3].

% Paper structure
The rest of this paper is organized as follows: Section \ref{sec:related} reviews related work.
Section \ref{sec:method} describes our methodology. Sections \ref{sec:experiments} and \ref{sec:results} present
our experimental setup and results. Section \ref{sec:discussion} discusses the implications,
and Section \ref{sec:conclusion} concludes with future work.

\section{Related Work}
\label{sec:related}

% Organize by theme, not chronologically

\subsection{Theme One}
Prior work on [theme one] has primarily focused on [main approaches].
Smith et al. \cite{smith2020} introduced [specific contribution], which [description].
However, this approach has limitations in [specific limitation].

Johnson and Lee \cite{johnson2019} extended this work by [extension], achieving [result].
Despite these improvements, [remaining limitation].

\subsection{Theme Two}
In the area of [theme two], recent work includes [approach] by Chen et al. \cite{chen2021}.
Their method [description], but [limitation].

\subsection{Comparison with Prior Work}
Unlike prior approaches that [what they do], our method [what we do differently].
Table \ref{tab:comparison} summarizes the key differences.

\begin{table*}[t]
\centering
\caption{Comparison of our approach with prior work.}
\label{tab:comparison}
\begin{tabular}{lcccc}
\toprule
& Smith et al. \cite{smith2020} & Johnson \cite{johnson2019} & Chen \cite{chen2021} & Ours \\
\midrule
Property 1 & Value & Value & Value & \textbf{Value} \\
Property 2 & Value & Value & \textbf{Value} & Value \\
Property 3 & \textbf{Value} & Value & Value & Value \\
Complexity & O(n$^2$) & O(n log n) & O(n$^2$) & \textbf{O(n)} \\
\bottomrule
\end{tabular}
\end{table*}

\section{Methodology}
\label{sec:method}

% High-level overview
Our approach, [method name], consists of three main components:
(1) [Component A], (2) [Component B], and (3) [Component C].
Figure \ref{fig:overview} illustrates the overall architecture.

\begin{figure*}[t]
\centering
\includegraphics[width=0.8\linewidth]{figures/overview.pdf}
\caption{Overview of our methodology.}
\label{fig:overview}
\end{figure*}

% Detailed description
\subsection{Component A}
\label{sec:component_a}

Component A addresses [specific problem]. The key insight is [insight].

\paragraph{Mathematical Formulation}
Formally, we define [definition]. Given [inputs], we compute [output] using:
\begin{equation}
\label{eq:formula}
f(x) = \sum_{i=1}^{n} w_i \cdot \phi_i(x) + b
\end{equation}
where $w_i$ are [description], $\phi_i$ are [description], and $b$ is [description].

\paragraph{Algorithm}
Algorithm \ref{alg:component_a} describes the computation process.

\begin{algorithm}[t]
\caption{Component A Algorithm}
\label{alg:component_a}
\begin{algorithmic}[1]
\STATE {\bf Input:} $X, Y, \epsilon$ 
\STATE {\bf Output:} $\theta$ 
\STATE Initialize $\theta$ randomly
\FOR{iteration = 1 to max_iterations}
    \STATE Compute gradient: $g \leftarrow \nabla_\theta J(\theta)$
    \STATE Update: $\theta \leftarrow \theta - \alpha \cdot g$
    \STATE Check convergence: if $|g| < \epsilon$, break
\ENDFOR
\RETURN $\theta$
\end{algorithmic}
\end{algorithm}

\subsection{Component B}
\label{sec:component_b}

Component B [description]. The implementation uses [technology/algorithm].

\paragraph{Complexity Analysis}
The time complexity of Component B is O([complexity]), and space complexity is O([complexity]).

\subsection{Component C}
\label{sec:component_c}

\paragraph{Implementation Details}
We implemented Component C using [language/framework].
Key optimizations include:
\begin{itemize}
    \item [Optimization 1]: [description]
    \item [Optimization 2]: [description]
    \item [Optimization 3]: [description]
\end{itemize}

% Theoretical properties
\subsection{Theoretical Properties}

\begin{theorem}
\label{thm:convergence}
Under [assumptions], Algorithm \ref{alg:component_a} converges to a local minimum at rate O($1/t$).
\end{theorem}

\begin{proof}
[Proof here]
\end{proof}

\input{sections/experiments}
\input{sections/results}
\input{sections/discussion}
\input{sections/conclusion}

% References
\bibliographystyle{IEEEtran}
\bibliography{references}

\end{document}
